\chapter*{Hướng Dẫn Sử Dụng Sách}
\addcontentsline{toc}{chapter}{Hướng Dẫn Sử Dụng Sách}
\markboth{Hướng Dẫn Sử Dụng Sách}{Hướng Dẫn Sử Dụng Sách}

\begin{truyen}{Bộ Sách Này Dành Cho Ai?}{Who Is This Book For?}
\chuhoa{Q}{uyển sách này dành cho người trẻ đang lớn, người trưởng thành đang bận, và cả những ai đã từng trải nhưng vẫn muốn sống sáng suốt hơn.}
\textit{(This book is for young people growing up, adults living busy lives, and also for those already experienced but still wanting to live more wisely.)}

\textbf{Giáo viên:} dùng mỗi truyện như một lời gợi mở cuối giờ hoặc đầu tiết để học sinh nghĩ sâu hơn về đời sống.
\textit{(\textbf{Teachers:} use each story as an opening or closing reflection so students think more deeply about life.)}

\textbf{Học sinh:} mỗi ngày đọc một truyện, không cần vội. Điều quan trọng là dừng lại ở phần câu hỏi và tự trả lời.
\textit{(\textbf{Students:} read one story a day without rushing. What matters is pausing at the reflection questions and answering them honestly.)}

\textbf{Người đi làm:} đọc chậm để thấy mình trong những lựa chọn, những tiếc nuối, và những điều vẫn còn kịp sửa.
\textit{(\textbf{Working adults:} read slowly to see yourself in the choices, regrets, and things that are still not too late to fix.)}
\end{truyen}

\begin{tcolorbox}[
  colback=agedgold!8,
  colframe=agedgold!70!black,
  coltitle=black,
  fonttitle=\bfseries,
  title={Giải Thích Các Ô Màu Trong Sách \quad\textit{(Color-Box Guide)}},
  arc=4pt, boxrule=1pt, breakable, enhanced
]
\smallskip
\textbf{\color{inkdark}Ô Xanh Khói --- Câu Chuyện Chính (\textit{Main Story})}\\
Phần truyện song ngữ Việt -- Anh. Đọc tiếng Việt để nắm ý, đọc thêm tiếng Anh để nhớ sâu hơn và luyện ngôn ngữ.\\
\textit{(The bilingual story section. Read Vietnamese for meaning, then the English to remember more deeply and practise language.)}

\medskip
\textbf{\color{smokeblue}Ô Đậm --- Điều Nên Giữ Lại (\textit{Key Lesson})}\\
Một ý cốt lõi sau mỗi truyện. Đây là câu nên gạch chân hoặc chép lại.\\
\textit{(The core lesson after each story. This is the line worth underlining or copying down.)}

\medskip
\textbf{\color{smokeblue!60}Ô Nhạt --- Easy English}\\
Hai câu tiếng Anh thật ngắn để đọc thành tiếng và ghi nhớ nhanh.\\
\textit{(Two very short English lines to read aloud and remember quickly.)}

\medskip
\textbf{\color{agedgold!80!black}Ô Vàng Đồng --- Gợi Ý Suy Ngẫm (\textit{Reflection})}\\
Hai câu hỏi để tự soi lại bản thân. Phần này không cần trả lời hay, chỉ cần trả lời thật.\\
\textit{(Two questions for self-reflection. You do not need perfect answers, only honest ones.)}
\smallskip
\end{tcolorbox}

\begin{baihoc}
Cách đọc tốt nhất là một truyện mỗi ngày. Cuốn sách này không cần đọc nhanh; nó cần được để lại trong đầu lâu hơn.
\textit{(The best way to read this book is one story a day. It is not meant to be read fast; it is meant to stay in your mind longer.)}
\end{baihoc}

\begin{ghinhoanh}
\textbf{Reading tips:}

Read slowly. Think honestly.

One story today can save one regret tomorrow.
\end{ghinhoanh}

\ngancach